\section{E3: a different objective for layer 1}
\label{sec:objective}

\lead{The prediction.} Classification rewards clauses that are specific, and specific clauses
fire rarely. Reconstruction has no such pressure --- a clause earns its keep by predicting
\emph{any} part of the input --- so an autoencoder-trained layer~1 should be denser, have
fewer dead channels, and preserve information a classifier discards. If the first report's
diagnosis was right, that should be enough to unblock layer~2.

\lead{The objective learns.} Predicting the \aeTarget\ centre bits of a patch from its
\aeContext\ surrounding ones, over \aePatches{}k patches and \aeEpochs\ epochs, reaches
\aeRecon{}\% per-bit reconstruction accuracy against a \aeBaseline{}\% majority-class
baseline. It is a real task and the clauses solve it.

\subsection{Result: the statistics change exactly as predicted}

\begin{table}[h]
\centering\small
\caption{Layer-1 statistics by objective, measured on 5000 training images.
\code{fire} is the median per-channel firing rate of the $\poolK\times\poolK$ OR-pooled map ---
what layer~2 actually sees --- \code{dead} the fraction of channels that never fire,
\code{size} the median clause size, and \code{neg frac} the mean fraction of each clause's
literals that are negations.}
\label{tab:l1stats}
\pgfplotstabletypeset[
  col sep=comma, string type,
  every head row/.style={before row=\ttoprule, after row=\ttmidrule},
  every last row/.style={after row=\ttbotrule},
  columns={label,fire,dead,size,neg_frac},
  columns/label/.style={column name=\thead{layer-1 objective}, column type={l}},
  columns/fire/.style={column name=\thead{fire (\%)}, column type={r}},
  columns/dead/.style={column name=\thead{dead (\%)}, column type={r}},
  columns/size/.style={column name=\thead{size}, column type={r}},
  columns/neg_frac/.style={column name=\thead{neg frac}, column type={r}},
]{data/l1_stats.csv}
\end{table}

The mechanism works. Every statistic the first report named as the problem moves in the right
direction, and by a large factor rather than a marginal one.

\subsection{And on CIFAR-10 the stack still falls short}

\begin{ttkey}
On CIFAR-10 the autoencoder stack reaches \accMctmAuto{}\,$\pm$\,\stdMctmAuto{}\% against
\accMctm{}\% for the greedy stack --- a real gain, and criterion~1 is met. It is not close to
criterion~2: \emph{random} three-literal clauses still reach \accMctmRandom{}\%, and
criterion~3, the single layer at the same budget, is at \accSingleMatched{}\%. A better
objective for layer~1 moves the stack a third of the way from the greedy result to the random
control, and stops.

Hold that verdict until \S\ref{sec:synthetic}. The same layer-1 objective, on the tasks that
actually require what the second layer supplies, is the strongest result in this report; and
\S\ref{sec:randcalib} finds that what \ideaA\ is doing to the first layer's statistics can be
done to an \emph{untrained} layer directly, at which point CIFAR-10 works too.
\end{ttkey}

\begin{figure}[h]
\centering
\begin{tikzpicture}
\begin{axis}[
  width=0.92\textwidth, height=6.4cm,
  xlabel={epoch}, ylabel={test accuracy (\%)},
  grid=major, grid style={ttline!40},
  legend pos=south east, legend columns=2,
  legend style={font=\sffamily\scriptsize, draw=ttline},
  label style={font=\sffamily\small}, tick label style={font=\sffamily\footnotesize},
  every axis plot/.append style={very thick},
]
\addplot[ttslate] table[col sep=comma, x=epoch, y=test_acc] {data/curve_single-matched.csv};
\addplot[ttorange, dashed] table[col sep=comma, x=epoch, y=test_acc] {data/curve_mctm-random.csv};
\addplot[ttmoss] table[col sep=comma, x=epoch, y=test_acc] {data/curve_mctm-auto.csv};
\addplot[ttamber] table[col sep=comma, x=epoch, y=test_acc] {data/curve_mctm.csv};
\legend{single matched, MCTM random L1, MCTM autoencoder L1, MCTM greedy L1}
\end{axis}
\end{tikzpicture}
\caption{Test accuracy per epoch (seed 0). The $x$ axis counts layer-2 epochs; layer~1 is
trained and frozen beforehand in every arm shown.}
\label{fig:curves-auto}
\end{figure}

\subsection{What the readout says, and what it cannot say}
\label{sec:readout}

The claim \ideaA\ rests on is about \emph{information}: an unsupervised layer~1 should keep
what a classifier throws away. The first report measured that with the \code{flat-stack}
readout --- globally OR-pool the layer-1 bits, train a flat machine on them --- and the same
readout run on the autoencoder layer gives \accFlatAuto{}\% against \accFlatStack{}\% for the
greedy layer and \accFlatRandom{}\% for random clauses.

That ordering is the reverse of the one \ideaA\ predicts, and the reason is in
Table~\ref{tab:l1stats} rather than in the features. A channel firing on $p$ of the 841 patch
positions is globally on with probability $\approx 1-(1-p)^{841}$; at $p$ of a few per cent
that is indistinguishable from~1. A layer dense enough to be useful to a \emph{spatial} second
layer is simultaneously saturated for a global-OR readout, so the readout measures density,
not information, once density is above a fraction of a percent. It is a valid comparison among
sparse layers --- which is how the first report used it --- and not a valid one here. The
honest measurement of \ideaA\ is the \mctm\ row, and it says the objective helped a little.
